% Figure: the six trigonometric ratios read as lengths on the unit circle.
% A radius at angle theta meets the circle at P = (cos, sin). The vertical
% dropped from P gives sin; the horizontal gives cos. The tangent line at
% the point (1,0) extended to the radius gives tan (vertical intercept) and
% sec (length along the radius). The line tangent at (0,1) gives cot and csc.
\begin{figure}[htbp]
\centering
\begin{tikzpicture}[scale=2.6, line cap=round, >=latex]
  % axes
  \draw[->, engblack, thin] (-0.25, 0) -- (1.75, 0) node[right, font=\small] {$x$};
  \draw[->, engblack, thin] (0, -0.25) -- (0, 1.75) node[above, font=\small] {$y$};
  % unit circle
  \draw[enggray, thick] (0,0) circle (1);
  % chosen angle theta = 50 degrees; cos ~ 0.6428, sin ~ 0.7660, tan ~ 1.1918, sec ~ 1.5557
  % radius to P, extended to the vertical tangent line x = 1 at height tan
  \draw[engblue, very thick, ->] (0,0) -- (1.5557, 1.1918);
  \filldraw[engblue] ({cos(50)}, {sin(50)}) circle (0.02);
  \node[engblue, font=\footnotesize, anchor=south east] at ({cos(50)}, {sin(50)}) {$P$};
  % vertical tangent line at x = 1
  \draw[engamber, thick] (1, 0) -- (1, 1.35);
  % sin drop from P
  \draw[engred, thick] ({cos(50)}, 0) -- ({cos(50)}, {sin(50)});
  \node[engred, font=\footnotesize, anchor=east] at ({cos(50)}, {0.5*sin(50)}) {$\sin\theta$};
  % cos along x-axis
  \draw[engred, thick] (0,0) -- ({cos(50)}, 0);
  \node[engred, font=\footnotesize, anchor=north] at ({0.5*cos(50)}, 0) {$\cos\theta$};
  % tan is the segment along x=1 from (1,0) to (1, tan)
  \draw[engamber, very thick] (1, 0) -- (1, 1.1918);
  \node[engamber, font=\footnotesize, anchor=west] at (1, 0.6) {$\tan\theta$};
  % sec is the length of the radius from origin to (1, tan)
  \node[engblue, font=\footnotesize, anchor=south east, rotate=37.5] at (0.78, 0.60) {$\sec\theta$};
  % angle arc
  \draw[engblack] (0.28, 0) arc (0:50:0.28);
  \node[engblack, font=\footnotesize] at (0.38, 0.16) {$\theta$};
  % right-angle mark at (1,0)
  \draw[engblack, thin] (0.93, 0) -- (0.93, 0.07) -- (1, 0.07);
\end{tikzpicture}
\caption[Trigonometric ratios as lengths]{The trigonometric ratios read as
lengths on the unit circle. A radius at angle $\theta$ meets the circle at
$P = (\cos\theta, \sin\theta)$, so the horizontal and vertical legs of the
inscribed right triangle are the cosine and the sine directly. Extending the
radius to the vertical line $x = 1$ (tangent to the circle at $(1,0)$) reads
the tangent as the intercept height $\tan\theta$ and the secant as the length
of the extended radius $\sec\theta$, because the small triangle on the
tangent line is similar to the inscribed one with its adjacent side scaled to
unity. The names \emph{tangent} and \emph{secant} record this construction: the
tangent is measured on the tangent line, and the secant is measured along the
line that cuts the circle. The picture makes $1 + \tan^{2}\theta =
\sec^{2}\theta$ the Pythagorean theorem of the tangent-line triangle.}
\label{fig:vol02ch03:secant-tangent-lengths}
\end{figure}
